\chapter{Introduction}
\label{ch:introduction}

\section{Background}

Consumer unmanned aerial vehicles (UAVs, or \emph{drones}) are cheap, widely available, and increasingly implicated in unauthorised surveillance, perimeter incursion, and disruption of civil airspace~\cite{shi2018counteruas,taha2019drone,samaras2019deep}. Protecting a fixed site against them imposes two demands at once. The system needs \emph{high recall}, because a missed drone defeats its purpose. It also needs a \emph{low false-alarm rate}, because constant alerts for birds and passing aircraft desensitise the operator it is meant to serve.

No single sensing modality meets both cheaply, and the Counter-UAS (C-UAS) literature explains why~\cite{samaras2019deep}. Radar struggles with the small radar cross-section of quadcopter-class drones. Radio-frequency detection needs an active control link, which fiber-optic-guided drones (fielded at scale since 2024~\cite{csis2025dronewar}) lack by construction, and which autonomous drones can forgo. Acoustic arrays degrade under wind, traffic, and aircraft noise. Visual detection with deep neural networks is passive, inexpensive, and reads the surrounding scene. It is the modality this thesis adopts.

Modern single-stage detectors are genuinely strong at this task. The production RGB detector built here reaches $P=0.989$, $R=0.982$ on the Anti-UAV benchmark~\cite{jiang2021antiuav}, with only 41 false positives across 4{,}000 frames, and it hallucinates on just $2.8\%$ of its in-distribution test dataset.\footnote{Anti-UAV frames appear in the detector's training dataset, so this is an in-distribution sanity floor rather than a generalisation claim (Section~\ref{sec:lit_drone}); evaluated at \texttt{imgsz=640}, $n=4{,}000$.} On ordinary footage (sky, buildings, traffic, vegetation), these detectors already solve the false-alarm problem at the detection head.
% [source: thesis_eval/results/tier1_results.json (antiuav A ft4/rgb); run=thesis_eval/pipeline_eval_unified.py; config=tier1 antiuav @640 n=4000]
% [source: eval=v5_rgbds (bare halluc 0.028); config=rgb_dataset_iou_640]

\section{Problem Statement}
\label{sec:problem}

The false-alarm problem is not generic. On most footage the detector never fires falsely at all; the failures concentrate in scenes with the three aerial \emph{confuser} classes: birds, fixed-wing aircraft, and helicopters. On a 2{,}633-image out-of-distribution confuser dataset assembled for this thesis, the same detector fires on $30.4\%$ of frames, and on bird-only Svanstr\"om footage that rises to $94\%$ (the per-category breakdown is in Table~\ref{tab:rgb_comparison}; example hallucinations in Figure~\ref{fig:confuser_fp_examples}).
% [source: thesis_eval/results/tier1_results.json (rgb_confuser C bare); thesis_eval/results/notes_round1_results.json (rgb_confuser CAT); config=tier1 rgb_confuser @640 n=2633]

The cause is geometric. At surveillance range, all four classes (drone, bird, aeroplane, helicopter) appear as small, low-texture silhouettes against sky, and the gap that separates them is too narrow for the detector to exploit \emph{at its decision layer} without losing drones. The separation does exist, but deeper in the network, inside the detector's own intermediate features (Section~\ref{sec:model_mri}). Exploiting it \emph{there}, rather than at the detection head, is the central idea of this thesis.
% [source: ledger=retrainedv2-recall-collapse; eval=rgb_svan_baseline; cache=eval/results/_failure_diagnosis/svanstrom_1280_by_category.csv; config=svan_iop_1280]

The obvious fix, retraining the detector with confuser images as hard negatives, has class-specific and diminishing returns. Training with explicit confuser negatives lowers helicopter and aeroplane fire but leaves bird fire essentially unchanged (Table~\ref{tab:rgb_comparison}). Training more aggressively does suppress birds, but collapses Svanstr\"om drone recall from $R=0.961$ to $R=0.306$ at the same precision. The collapse is specific to out-of-distribution footage: on its own in-domain test split the same weights score $F1=0.949$, \emph{above} the baseline's $0.942$. In-domain model selection would therefore pick exactly the detector that fails in the field. And recall, once lost at the detector, cannot be recovered downstream: no later stage can add a detection the detector never proposed.
% [source: ledger=retrainedv2-recall-collapse; eval=rgb_svan_hardneg,rgb_svan_retrainedv2; cache=eval/results/_failure_diagnosis/svanstrom_1280_by_category.csv; config=svan_iop_1280]
% [source: ledger=rgb-collapse-ood-specific; eval=rgb_rgbds_retrainedv2 (0.9494 in-domain > baseline 0.942); config=rgb_dataset_iou_640]

Two further gaps separate benchmark scores from field readiness. The first is \emph{resolution}: small or distant drones fall below the detector's resolvable floor and disappear entirely (Section~\ref{sec:resolution_arch}). On Svanstr\"om, where the median drone is only $29.8$~px across, bare recall falls from $0.90$--$0.96$ in the $16$--$64$~px range to $0.63$ below $16$~px (Section~\ref{sec:per_size}). The second is \emph{modality}: the visible channel degrades at night and under glare, exactly where the thermal channel is strongest, so the system is dual-modality from the start. Which modality wins \emph{reverses} with the scene: thermal leads on Svanstr\"om ($F1$ $0.940$ vs $0.607$), visible leads on Anti-UAV ($0.985$ vs $0.961$), and the routed system matches whichever is stronger (Table~\ref{tab:rq3}).
% [source: thesis_eval/results/notes_round1_results.json (svanstrom SZ, M); appendix tab Svanström native 640x512]

Any discrimination the detector cannot learn without costing recall must therefore move \emph{off} the detector. This thesis relocates confuser rejection, modality choice, and false-positive filtering into a downstream system: a \emph{trust router} that picks the reliable modality each frame, a per-detection \emph{confuser filter} that re-reads the detector's own internal features to remove false positives, and a \emph{temporal smoother} that turns per-frame decisions into alerts (Section~\ref{sec:distill_verifier}). The detectors themselves are left tuned purely for recall.

One evaluation surface comes from a deployment partner: a fixed perimeter-CCTV camera (SelCom) whose 1080p stream renders drones at a median of $\approx$37~px amid urban clutter. The system is not deployed there, but the partner's footage supplied a fine-tuning and test surface. Fine-tuning lifted its validation $F1$ from $0.145$ to $0.591$, while leaving the in-domain test dataset essentially unchanged ($0.942 \to 0.929$). Two engineering constraints hold throughout: every component must run interactively on a single consumer GPU (Section~\ref{sec:pipeline_speed}), and every reported number must come from a configuration the system can actually ship.
% [source: dataset=selcom_cctv; file=SelCom CCTV manifest (median 36.8px native)]
% [source: eval=v5_selcom_bare (ft4 0.5911), selcom baseline 0.145 (eval/results _selcom cmp); eval=rgb_rgbds_retrainedv2 extra (baseline 0.942, ft4 0.929)]

\section{Research Questions}
\label{sec:rqs}

The thesis is organised around three research questions; the answers are interpreted in Chapter~\ref{ch:discussion}.

\begin{description}
    \item[RQ1 (capability)] Retraining the detector was the cheap first answer, and it failed (Section~\ref{sec:problem}): hard-negative training either leaves the dominant confuser class untouched or collapses small-drone recall. \emph{Where retraining fails}, can a layered pipeline maintain or improve drone detection while suppressing false positives, on both ordinary and confuser-rich scenes?
    \item[RQ2 (attribution)] What does each stage contribute? Measured by ablation (the trust router, the confuser filter, their composition order, and temporal aggregation), each isolated on the same frames, with the runtime cost of each stage accounted for.
    \item[RQ3 (dual-modality value)] Does per-frame routing between two modalities beat either modality alone? The scoring protocol this question requires is defined in Section~\ref{sec:metrics}.
\end{description}

A single methodological discipline underpins the downstream stages: \emph{measure before you train}. Before training the router or the filters, the Model MRI (Section~\ref{sec:model_mri}) measured what the detectors' features can already separate, and a leakage statistic excluded features that merely fingerprint the scene. The detectors themselves were trained conventionally; separately, a human-in-the-loop process co-developed the thermal detector with its data.

\section{Contributions}

This thesis makes four contributions and reports one empirical finding:

\begin{enumerate}
    \item \textbf{A deployable dual-modality drone-detection system with an operator GUI.} An RGB and a thermal detector are combined by a learned \emph{trust router} that chooses, per frame, which modality to believe; a per-detection \emph{confuser filter} removes false positives; and a temporal smoother converts per-frame decisions into alerts. Held to one declared evaluation standard (Section~\ref{sec:eval_protocol}), the full pipeline lifts Svanstr\"om drone $F1$ from $0.742$ to $0.946$ \emph{while raising recall}, does no harm on the saturated Anti-UAV control, and cuts out-of-distribution confuser fire by more than an order of magnitude, with every stage running on every frame within budget (Chapter~\ref{ch:hitl}).
    % [source: thesis_eval/results/tier1_results.json (svanstrom B, rgb_confuser C, antiuav B); run=thesis_eval/pipeline_eval_unified.py]
    % [source: ledger=robust6-speed-feature-efficiency,v5-ship-per-frame; run=eval/bench_speed.py]
    \item \textbf{The Model MRI, a statistics-before-training instrument for detector feature spaces.} It reads a detector's own internal features for each detection and answers design questions with classical statistics \emph{before} any model is trained, chiefly how separable two classes (drone vs.\ confuser) already are. Where they are separable, it trains the confuser filter directly on those same features (drone detections as positives, confuser detections as negatives). Measured drone/confuser separability was high in both modalities ($0.952$ RGB, $0.981$ IR). These in-dataset diagnoses license a training attempt; the out-of-distribution evaluations of Chapter~\ref{ch:hitl} are what validate it.
    % [source: ledger=v5-lda-separability,ir-features-separable; run=mri/cli.py]
    \item \textbf{A human-in-the-loop data engine.} A purpose-built label reviewer and a disciplined dataset--model co-evolution loop, reported as a six-revision case study of the thermal detector. Scored throughout on one fixed held-out test split, the detector improved from $F1$ $0.503$ to $0.967$. The trace includes the one revision where the review step was skipped (V5): an unreviewed batch carried unlabelled drones, cost $12.7$ points of precision, and the same review loop is what later found and corrected it.
    % [source: ledger=ir-version-progression; eval=ir_final_* (all on config ir_final_640, same split)]
    \item \textbf{A set of curated data artifacts.} The product of the manual review behind the system: a $129{,}130$-frame reviewed thermal dataset; a $27{,}024$-frame RGB confuser dataset with sequence-disjoint splits; the $28{,}710$-frame paired Svanstr\"om conversion to YOLO format; and a $19$-clip real-video benchmark ($2{,}609$ frames, $1{,}234$ annotated drone instances) that is fully out-of-distribution for every component. Licences preclude redistributing the source datasets, so these are claimed as reviewed artifacts; the real-video benchmark's labels are publishable (Appendix~\ref{app:datasets}).
    % [source: tab:ds_ir_components total 129,130; rgb_confusers_merged 27,024; svanstrom_paired 28,710; sec:ds_youtube 19 clips / 2,609 frames / 1,234 GT instances]
    \item \textbf{An unplanned cross-modal finding: thermal-to-grayscale transfer.} Run zero-shot on grayscale-converted RGB, the thermal-trained detector still detects drones when no thermal camera is available, landing within a few points of the dedicated RGB detector on the same frames ($F1$ $0.580$ vs $0.607$ on Svanstr\"om) and tying it on the hardest bird-cluttered clip ($0.837$ vs $0.840$). A raw-RGB control isolates the grayscale conversion itself as the mechanism, and we are not aware of this transfer being reported for this thermal-to-grayscale setup; it gives the system a usable RGB-only fallback (Section~\ref{sec:grayscale}).
    % [source: thesis_eval/results/tier1_results.json (D table); ledger=ir-grayscale-fallback (clip + 0.837/0.840); eval=patch_catch_v2_svan,v5_svan_mlp/patch]
\end{enumerate}

Every number in this thesis is held to one standard: a single canonical configuration per dataset, per-modality trust-aware scoring, frame counts printed alongside, and bootstrap confidence intervals on headline results (Section~\ref{sec:eval_protocol}). Every reported number is also traceable: each is a named cell in a cached results file, regenerable by a one-line command and checked by an automated audit (Appendix~\ref{app:provenance}).

\section{Ethical Considerations and Dual-Use}
\label{sec:ethics}

Counter-UAS sensing is dual-use technology. The same passive visual-detection capability that flags a hostile drone over critical infrastructure could, in a different deployment, contribute to mass aerial surveillance or armed-response targeting. This thesis develops the detection component only: it neither integrates with effectors (jammers, kinetic intercept, take-down systems) nor classifies operator behaviour. The training datasets contain no personally identifying information, no identifiable third-party imagery beyond what is already public in the cited datasets and clips, and no acoustic or RF data that could identify an operator. The deployment scope assumed throughout is the protection of fixed sites (perimeters, airport approaches, industrial facilities), where a detected drone is reported to a human operator who decides the response. Because every result is reproducible from the published caches (Appendix~\ref{app:provenance}), any future adaptation of the system remains auditable.

\section{Thesis Outline}
\label{sec:outline}

Chapter~\ref{ch:literature} reviews the literature: real-time object detection, drone detection and the confuser problem, thermal imaging, multi-modal fusion, cross-modal transfer and cascaded rejection, and the two methodological lineages this work builds on (data-centric annotation and feature-space probing). Chapter~\ref{ch:architecture} presents the system: the pipeline and each component, with the design rationale. Chapter~\ref{ch:methodology} gives the experimental method: the datasets and their roles, the Model MRI and the human-in-the-loop engine, and the evaluation protocol. Chapter~\ref{ch:hitl} reports the results: the full-pipeline ablation that answers RQ1--RQ3, the per-component evidence (including the thermal detector's six-revision evolution), the Model MRI findings, and the grayscale cross-modal finding. Chapter~\ref{ch:discussion} interprets the answers to the research questions and sets out the threats, limitations, and future work. Chapter~\ref{ch:conclusion} states the production stack with its honest carve-outs and closes the thesis. The appendices detail the datasets (Appendix~\ref{app:datasets}), the models (Appendix~\ref{app:models}), and the glossary (Appendix~\ref{app:glossary}).
